% MT009 / explanation — canonical editable content.
数据有\(N=2000\)条、每条\(d=100{,}000\)个特征。线性SVM最后仍用\(f(x)=w^Tx+b\)预测，但求出这个模型可以走两条路线。


\textbf{原始问题（primal）按特征求权重。} \(w\)有100,000个分量，分别决定每一维特征怎样影响分数。这里优化变量规模主要跟特征数\(d\)走。


\textbf{对偶问题（dual）按训练样本求系数。} 给每条样本一个乘子\(\alpha_i\)，再用



\[
w=\sum_{i=1}^{N}\alpha_i y_i x_i
\]


把样本组合成同一个权重向量。\(y_i\)是正负标签，\(\alpha_i\)控制第\(i\)条样本的贡献；为0的项不参与最终预测。此时主要优化2000个样本系数，而非100,000个特征权重，所以\(d\gg N\)时可以考虑对偶求解器。


对偶还要计算样本之间的内积。把\(K_{ij}=x_i^Tx_j\)排成表就是Gram矩阵，它有\(2000\times2000=4{,}000{,}000\)个元素，并非\(100{,}000^2\)。构造这些内积本身仍有成本，减少未知数不自动等于每种实现都更快。


实际建议是比较合适的primal/dual线性SVM求解器，并利用稀疏输入：若很多特征为0，不要按稠密矩阵反复运算。也可验证降维或特征筛选，但需检查是否损失分类信息。英文答题重点给出\(d\)、\(N\)的比较及换用对偶的理由，不必完整推导拉格朗日对偶。


\textbf{出处：} 2020B期中Quiz Q9，原卷第4页（10分）


\textbf{答案依据：} 2020B期中 Q9 随卷／配套参考答案；条件与解释经复核。
